\begin{proof}[Proof of \cref{obj:lem:signed-depth-stationary-law}]
Write
\[
\bar d_r(0,(j,u))
:=
w_j\,\frac{1+u\varepsilon_C r^j}{2},
\qquad
j\in\{0,\ldots,Q\},\quad u\in\{-1,+1\}.
\]
We prove that this probability vector is stationary for \(P^{(r)}\), and then use the
contraction statement to identify it with the selected stationary law.

\begin{enumerate}
\item Since \(t_0>0\), the value \(\alpha=e^{-1/t_0}\) lies in \((0,1)\). Also
\(C\ge 1\) gives \(0\le \varepsilon_C\le 1/2\), and \(r\in[0,1]\) gives
\(0\le r^j\le 1\). Hence every factor
\((1+u\varepsilon_C r^j)/2\) is nonnegative. For each fixed depth \(j\),
\[
\sum_{u\in\{-1,+1\}}
w_j\,\frac{1+u\varepsilon_C r^j}{2}
=w_j,
\]
and the geometric normalization gives
\[
\sum_{j=0}^Q w_j
=
\frac{1-\alpha}{1-\alpha^{Q+1}}\sum_{j=0}^Q \alpha^j
=1.
\]
Thus \(\bar d_r\) is a probability vector on \(\mathcal S_Q\). % lean: CausalSmith.Stat.PomdpLatentOverlapMinimax.signedDepthStationaryWeight_probabilityVector

\item The state kernel induced by the constant policy \(a_r\) has the following
hidden-state transition probabilities, obtained from the signed-depth dynamics in
\cref{obj:def:signed-depth-family} by averaging the action-one and action-zero
advance weights. If \(j<Q\), then
\[
P^{(r)}\bigl((0,(j,u)),(0,(0,u'))\bigr)
=(1-\alpha)\frac{1+u'\varepsilon_C}{2},
\]
and
\[
P^{(r)}\bigl((0,(j,u)),(0,(j+1,u'))\bigr)
=
\alpha\,\frac{1+uu'r}{2}.
\]
If \(j=Q\), then
\[
P^{(r)}\bigl((0,(Q,u)),(0,(0,u'))\bigr)
=
\frac{1+u'\varepsilon_C}{2}.
\]
The remaining transitions have weight zero. % lean: CausalSmith.Stat.PomdpLatentOverlapMinimax.signedDepth_policyKernel_apply

For a target state of depth \(0\), summing the displayed reset probabilities gives
\[
(\bar d_rP^{(r)})(0,(0,u'))
=
\frac{1+u'\varepsilon_C}{2}
\left\{
w_Q+(1-\alpha)\sum_{j=0}^{Q-1}w_j
\right\}.
\]
Since \(\sum_{j=0}^Q w_j=1\),
\[
w_Q+(1-\alpha)\sum_{j=0}^{Q-1}w_j
=
(1-\alpha)+\alpha w_Q
=
(1-\alpha)+
\frac{(1-\alpha)\alpha^{Q+1}}{1-\alpha^{Q+1}}
=
\frac{1-\alpha}{1-\alpha^{Q+1}}
=w_0.
\]
Therefore
\[
(\bar d_rP^{(r)})(0,(0,u'))
=
w_0\,\frac{1+u'\varepsilon_C}{2}
=
\bar d_r(0,(0,u')).
\]

For a target state of depth \(m\in\{1,\ldots,Q\}\), only depth \(m-1\) can advance
to it, so
\[
(\bar d_rP^{(r)})(0,(m,u'))
=
\alpha w_{m-1}
\sum_{u\in\{-1,+1\}}
\frac{1+u\varepsilon_C r^{m-1}}{2}
\frac{1+uu'r}{2}.
\]
The sign sum is
\[
\sum_{u\in\{-1,+1\}}
\frac{1+u\varepsilon_C r^{m-1}}{2}
\frac{1+uu'r}{2}
=
\frac{1+u'\varepsilon_C r^m}{2}.
\]
Using \(w_m=\alpha w_{m-1}\), this yields
\[
(\bar d_rP^{(r)})(0,(m,u'))
=
w_m\,\frac{1+u'\varepsilon_C r^m}{2}
=
\bar d_r(0,(m,u')).
\]
Thus \(\bar d_rP^{(r)}=\bar d_r\). % lean: CausalSmith.Stat.PomdpLatentOverlapMinimax.signedDepthStationaryWeight_isStationary

\item The contraction supplied by \cref{obj:lem:signed-depth-constant-policy-contraction}
implies uniqueness of stationary probability vectors. Indeed, if \(\lambda\) and
\(\lambda'\) are stationary for \(P^{(r)}\), then
\[
\|\lambda-\lambda'\|_{\mathrm{TV}}
=
\|\lambda P^{(r)}-\lambda'P^{(r)}\|_{\mathrm{TV}}
\le
\alpha\|\lambda-\lambda'\|_{\mathrm{TV}}.
\]
Because \(\alpha<1\) and total variation is nonnegative, the norm is zero, hence
\(\lambda=\lambda'\) on the finite state space. The selected law \(d_r\) is
stationary because the stationary vector \(\bar d_r\) exists, so uniqueness gives
\(d_r=\bar d_r\). % lean: CausalSmith.Stat.PomdpLatentOverlapMinimax.stationary_unique_of_contraction

\item Evaluating the identity \(d_r=\bar d_r\) at the hidden state
\(h=(j(h),u(h))\) gives
\[
d_r(0,h)
=
w_{j(h)}\,\frac{1+u(h)\varepsilon_C r^{j(h)}}{2},
\]
as claimed. % lean: CausalSmith.Stat.PomdpLatentOverlapMinimax.signedDepth_stationaryLaw_apply
\end{enumerate}
\end{proof}
